% ========================================
% Document Class and Package Setup
% ========================================
\documentclass[%
 reprint,
 amsmath,amssymb,
 aps,
]{revtex4-1}

\usepackage{graphicx}% Include figure files
\usepackage{dcolumn}% Align table columns on decimal point
\usepackage{bm}% bold math
\usepackage[dvipdfm,colorlinks=true,linkcolor=blue,filecolor=blue, urlcolor=blue]{hyperref}
\usepackage{braket}
\usepackage{caption}
\usepackage{threeparttable}
\numberwithin{equation}{section}% (p.37)


% ========================================
% Document Begin
% ========================================
\begin{document}

    \captionsetup[figure]{justification=raggedright,
        labelfont={bf},name={Fig.},labelsep=period}
    \captionsetup[table]{
        labelfont={bf},name={Table.},labelsep=period}


% ========================================
% Title, Author, and Date
% ========================================
    \title{Perfect Contrast cannot be obtained in the Electron Double-Slit Experiment}
        \author{Satoshi Hanamura}
        \email{hana.tensor@gmail.com}
        \date{July 28, 2019}
% ========================================
% Abstract
% ========================================
        \begin{abstract}
            Conventionally, the wave of particles which through the double-slit is assumed plane waves. In this research, we considered that the interference fringes built up through the double-slit have a difference amplitudes between the case of electrons and the case of photons. The difference between the two fringes is in the troughs of the waves. In this research, it is hypothesized that the amplitudes of waves passing through the left and right slits are not even in the double-slit experiment of electrons. Computer simulations performed to obtain the results supporting this hypothesis. The concept that waves of different amplitudes pass through a double-slit is reasonably to have the notion that two spinor particles pass through each slit.
            
        \end{abstract}
    \maketitle


% = = * * = =  * * = = * * = = * * = = * * = = * * = = * * = = * *
% ========================================
% Section 1: INTRODUCTION
% ========================================
\section{\label{sec:INTRODUCTION}INTRODUCTION}


    In the Tonomura experiment \cite{Tonomura}, the dark part of the interference fringes is not clear. Tonomura immersed himself in laboratory experiments to clarify those dark lines. However, no matter how ingenious the dark part became clear, he came to a certain consideration. That is, in principle, it is impossible to obtain a dark part that is clear in the experiment.
    
    The point that two kinds of particles, -- photons and electrons --, make interference fringes by double-slit experiment in common. However, in the valley part of the interference fringe, the photon clearly becomes a dark part, whereas in the case of the electron, a clear valley can not be emerged (Fig. \ref{f_Hitachi_box=0}). The difference of the results occurred by the visibility are zero or non-zero. In the experiment of photons, the crest and the trough can be clearly distinguished, while that of the electron is ambiguous.

    Till date, the wave generated by the particle is regarded as a plane wave until the particles reached the double-slit. This study shall change this regarding.

    In this research, we focus on the difference between visibility and attempt to verify whether these differences can be explained by applying the previous electronic model \cite{Hanamura}. 

        \begin{figure}[h!]
            \centering 
            \includegraphics[scale=0.205]{f_Hitachi_box.jpg}
                \caption{There is no clear dark line on troughs in visible. Tonomura noted that it was hard to identify the deepest trough lines strictly because the visibility is non-zero. He had been looking for the reason why this phenomenon occurs for the later years of his life. Note: The author of this paper added the blue line boxes and created partial enlarged image from the original Hitachi's image.}
            \label{f_Hitachi_box=0}
        \end{figure}

% = = * * = =  * * = = * * = = * * = = * * = = * * = = * * = = * * 
% ========================================
% Section 2: MOTIVATION
% ========================================
\section{\label{sec:MOTIVATION}MOTIVATION}

    Let us define the expression of the dark part of the interference fringe clearly using mathematical formulas. Tonomura, who demonstrated single-electron buildup of an interference pattern successfully, sent an email to a professor on December 13, 2010. Tonomura noted that:

    \begin{quote}
        "$p(x)$ should have yielded value zero on the troughs of the interference fringes with the deflection of electron waves by biprism. However, experimental results; $p_{min}=0$, was hardly obtained. Is it possible to consider that the value of visibility could not achieve 100 percent in principle as far as by using the electron biprism, or furthermore, by using the electron beam with the charge?" \cite{Saiensu-sha} (note: translated in English by the author of this paper.) 
    \end{quote}

    Where $p(x)$ is a probability density that electrons appear at space coordinates $x$ on the screen.

    $P_{vis}$ is a visibility of the biprism interference pattern that is defined by $p_{max}$ which is the maximum value of $p(x)$ and by $p_{min}$ which is the local minimum value of $p(x)$.
        
        \begin{equation}
            (\mathrm{visibility}):P_{vis} \equiv \frac{p_{max}-p_{min}}{p_{max}+p_{min}}
            \ ,
        \end{equation}

        \begin{figure}[h!]
            \centering 
            \includegraphics[scale=1.0]{f_Px=0.png}
                \caption{The result of the interference fringes due to the difference in the value of $P_{vis}$. (a) $P_{vis}=1$: Crests and troughs could be clearly distinguished. This is the same as the interference fringe obtained by making a plane wave incident on the double slit. (b) $P_{vis} \neq 1$: Troughs are ambiguous because $p_{min}$ is non-zero. For example, inside the circle surrounded by a yellow dotted line. Electrons dotted slightly even at the deepest trough part of the interference fringes.}
            \label{f_Px=0}
        \end{figure}


    \noindent
    when
        
        \begin{equation}
            p_{min}=0 \ ,
        \end{equation}

    \noindent
    $p(x)$ equal to 1. $p_{min}=0$ means that electrons should not be observed at each deepest trough of the interference fringes. Let us see the two blue flames in Fig. \ref{f_Hitachi_box=0}. We cannot confirm complete dark lines between one crest and the next. This means $p_{min} \neq 0$.

    Tonomura had attempted to obtain $P_{vis}=1$ in his experiments over and over because of pursuing for his cutting-edge projects. However, the observed visibility did not achieve $P_{vis}=1$. After elaborated his trials, the result $P_{vis} \neq 1$ led him to the viewing that the results were not caused by the accuracy of the way of experiments, but by a theoretical reason that is fundamentally undiscovered.
    
    When plane wave passing through the biprism, it has been considered that the plane waves of electrons pass through, then interference fringes appeared consequently. Tonomura had some ideas that the reason why the visibility is not completely zero. One example idea is that the wave is not plane. We should consider cylinder wave instead of the plane wave\cite{Saiensu-sha}.

    Similar phenomena can be found in experiments using C$_{60}$ \cite{Arndt} and experiments using helium atoms \cite{Carnal}. From the facts of such experimental results, the interference fringes of electrons emitted from the electron beam are to be searched for the reason in principle.


% = = * * = =  * * = = * * = = * * = = * * = = * * = = * * = = * * 
% ========================================
% Section 3: Applying the Model
% ========================================
\section{\label{sec:Applying_the_Model}Applying the Model}


% ----------------------------------------
% Subsection 3.1: Review the Previous Study
% ----------------------------------------
\subsection{\label{sub:Review_the_Previous_Study}Review the Previous Study}


    In this section, we shall apply the electron model \cite{Hanamura} to the two-slit experiment for obtaining a foundation of achieving the observed visibility $P_{vis}(x) \neq 1$.

    Let us glance at the points of the previous paper here. We assumed the images about an electron as follows:

    \begin{itemize}
        \item Assumption 1: An electron has an internal structure, which acquired three oscillators composed of one vector oscillator as a virtual photon and two spinor oscillators as bare electrons.
        \item Assumption 2: The bare electron is a thermal spot and a perfect black body. It radiates its energy harmonically through time-dependent oscillation.
        \item Assumption 3: While the two black bodies emit and absorb thermal energy alternately depending on the phase.
        \item Assumption 4: Replaced the image of a virtual photon with a real one. The virtual photon as a real photon captured by a bare electron with a force represented by the coupling constant of $\alpha \ (\approx 1/137)$.
        \item Assumption 5: The virtual photon can be moved as a simple harmonic oscillator with the emergence and disappearance of bare electrons at fixed spatial points x = a and x = --a. The two bare electrons would not change their respective spatial positions, and only the thermal energies of both the bare electrons are observed to change with time.
    \end{itemize}

    As a result of incorporating the function of zero point energy into the law of conservation of energy with the classical oscillator model, the three oscillators we obtained are as follows:


        \begin{equation}
            \begin{split}
                (Spinor 1): T_{\mathrm{e1}} \equiv E_{0} \cos^{4}{\left( \frac{\omega t}{2} \right)}, \\
                (Spinor 2): T_{\mathrm{e2}} \equiv E_{0} \sin^{4}{\left( \frac{\omega t}{2} \right)}, \\
                (Photon): \gamma_{\mathrm{K.E.}} \equiv \frac{1}{2} E_{0} \sin^{2}{(\omega t)},
                \end{split}
            \label{e_S1_S2_P}
        \end{equation}


    \noindent
    where $\omega t$ as the electron's phase. $T_{e1}$ and $T_{e2}$ are thermal potential energy (TPE) of the bare electrons, and $\mathrm{\gamma}^{*}_\mathrm{K.E.}$ is the kinetic energy of the virtual photon. The oscillator 1 would be a vector particle because it has $\omega t$ phase. The oscillator 2 and 3 would be spinor particles because they have $\omega t/2$ phase.

    These three oscillators have its own energy , $\mathrm{\gamma}^{*}_\mathrm{K.E.}$, $T_{e1}$ and $T_{e2}$, and total energy of these three corresponds to one free electron, $E_{0}$, as shown in eqn \ref{e_E=E1+E2+gamma}.

        \begin{equation}
            \begin{split}
                E_{0}
                    &= \ E_{0} \left( \mathrm{cos}^{4} \left( \frac{\omega t}{2} \right) + \mathrm{sin}^{4} \left( \frac{\omega t}{2} \right) + \frac{1}{2} \mathrm{sin}^{2} (\omega t) \right)
                    ,
            \end{split}
            \label{e_E=E1+E2+gamma}
        \end{equation}
    
    The first and the second term on the right side represents the spinors 1 and 2, and the third term represent the virtual photon. Because $\omega t$ implies as vector particle and $\omega t / 2$ would performed their phase as spinor particles, 720-degree rotation.

        \begin{figure}[h!]
            \centering 
            \includegraphics[scale=1.00]{f_QED.png}
                \caption{($\bm{\mathrm{a,b}}$) Self-energy diagram in QED and that in the modified model. Both diagrams $\bm{\mathrm{a}}$ and $\bm{\mathrm{b}}$ show the same example of “an electron” \cite{Hanamura}}.
            \label{f_QED}
        \end{figure}


% ----------------------------------------
% Subsection 3.2: Modification of the plane wave
% ----------------------------------------
\subsection{\label{sub:Modification_of_the_Plane_Wave}Modification of the plane wave}


    Electrons handled in the double slit experiment obtained high energy comparing to the ground state energy eqn \ref{e_E=E1+E2+gamma}.
    
    According to the previous electron model, one virtual photon is included in the electron in the ground state. Let us consider an electron in which a plurality of virtual photons are adsorbed to the electron. This extended model permit that when two spinors pass through the double slit with different routs, two spinors are wearing virtual photons in proportion to the TPE energy of each electron.

    Since the TPE taken by each spinor passing through the left and right changes periodically, the number and energy value of virtual photons surrounding the spinor also change accordingly. As this passes through the double slits, it is based on the mechanism that each amplitude is different in the circular wave generated from each slit.

    Presently, it is considered that plane waves pass in double slit experiments. Here we extend that concept. In the previous study \cite{Hanamura}, we examined the electron model in which two bare electrons are included in one electron. These bare electrons emit and absorb thermal potential energy (TPE) alternately depending on the phase, and their energy values fluctuate.
    
    The plane wave approximation passing through the double slit is not sufficient in this study. This is because, plane waves passing through the two slits are equal intensities. Let us consider a model that interferes waves with each of the two intensity values in which passes through a double slit.

        \begin{figure}[h!]
            \centering 
            \includegraphics[scale=1.00]{f_Biprism_1.png}
                \caption{($\bm{\mathrm{a}}$) An electron passes through the biprism. Picked up the two arrows drawn in Fig. \ref{f_QED}. ($\bm{\mathrm{b}}$) The bare electrons can be in various phase, on which are passing through biprism.}
            \label{f_Biprism_1}
        \end{figure}

    We extend the electronic model used in the previous study to explain how to occur the double-slit interference fringes. Fig. \ref{f_QED} was a modification of the Feynman diagram \cite{Hanamura}. In the modified electronic model, it can be reasonably understood that two bare electrons move in the direction of travel towards the slits. Electrons emitted from the electron gun split in the direction of travel before arriving at the Biprism. In this modified electron model, we consider that these bare electrons pass through the each slit as shown in Fig. \ref{f_Biprism_1} $(\bm{\mathrm{a}}, \bm{\mathrm{b}})$.

    
% ----------------------------------------
% Subsection 3.3: Modification of a fringe formula
% ----------------------------------------
\subsection{\label{sub:Modification_of_a_fringe formula}Modification of a fringe formula}


    The approximate expression for finding the interference fringes can be expressed by a Bessel function. The equation for interference fringes is to be improved so that the plane wave passing through the double slit can take two different amplitude values. The interference fringes can be expressed approximately with two circular waves using the zero-order Bessel function, $J_{0}$, as follows:

    \begin{equation}
        | \langle x_{i} , \psi \rangle | ^{2} = \left( J_{0}(k r_{1}) + J_{0}(k r_{2}) \right) ^{2}
        \label{e3_J0}
    \end{equation}

    Since we are based on the premise that each of the intensity of the electron wave when passing through the slit is different. In order to generate the wave, we shall add the weighting coefficients, $A_{\mathrm{spinor1}}$ and $A_{\mathrm{spinor2}}$, to the above wave equation;
    
    \begin{equation}
        | \psi (x) | ^{2} = \left( A_{\mathrm{spinor1}} \cdot J_{0}(k r_{1}) + A_{\mathrm{spinor2}} \cdot J_{0}(k r_{2}) \right) ^{2}
        \label{e3_J0_W}
    \end{equation}

        \begin{figure}[h!]
            \centering 
            \includegraphics[scale=1.00]{f_Biprism_2.png}
                \caption{($\bm{\mathrm{a}}$) Transplant the colored arrows to the $\theta$-axis. The transparency of the arrow and its phase are correspond to each other. Opaquer colored both blue and green arrows represent that the bare electrons have much thermal potential energy (TPE). $T_{\mathrm{e1}}$ and $T_{\mathrm{e2}}$ represent value of TPE of the two thermal spots as two spinors \cite{Hanamura}. ($\bm{\mathrm{b}}$) Schematic diagram in phase difference of the bare electrons when passing through the biprism. The red numbers are proportions of the amplitude of the bare electrons as they pass through the biprism.}
            \label{f_Biprism_2}
        \end{figure}

    We shall adopt,

        \begin{equation}
            \begin{split}
                A_{\mathrm{spinor1}} = A \mathrm{cos}^{2} \ \theta \ , \\
                A_{\mathrm{spinor2}} = A \mathrm{sin}^{2} \ \theta
                \ ,
            \label{e3_Amplitude_AB}
            \end{split}
        \end{equation}

    \noindent
    as weighting coefficients. Where $\theta$ denotes the electron's phase $\omega t$. The basis of these weighted values are that the energy value which can be taken by the bare electrons is from 0 to a positive value. It is because TPE does not take a negative value. Substitute these values for eqn \ref{e3_J0_W} and 1 for $A$ for simplicity;

    \begin{equation}
        | \psi (x) | ^{2} = \left( \mathrm{cos}^{2} \ \theta \cdot J_{0}(k r_{1}) + \mathrm{sin}^{2} \ \theta \cdot J_{0}(k r_{2}) \right) ^{2}
        \label{e3_psi_1}
    \end{equation}

    The wave function to be found is taking the integral of the above equation for $\theta$,

    \begin{equation}
        | \psi (x) | ^{2} = \int_{0}^{2 \pi} \left( \mathrm{cos}^{2} \theta \cdot J_{0}(k r_{1}) + \mathrm{sin}^{2} \theta \cdot J_{0}(k r_{2}) \right)^{2} d \theta
        \ .
        \label{e3_psi_2}
    \end{equation}


        \begin{figure}[h!]
            \centering 
            \includegraphics[scale=0.4]{f_6_edit.png}
                \caption{State of probability density of interference fringes changing with phase of an electron. Of the six diagrams, only the top 50 : 50 figure has a visibility value of 1. In the bottom diagram 100 : 0 shows $P_{vis} = 0$, therefore no interference fringes are observed.}
            \label{f_6_edit}
        \end{figure}
    

    Find all probability density functions across the $x$-axis of space, so the right side of eqn \ref{e3_psi_2} shall be integrated further in space;
    
    \begin{equation}
        \frac{1}{N} \int_{-\infty}^{\infty} \int_{0}^{2 \pi} \left( \mathrm{cos}^{2} \theta \cdot J_{0}(k r_{1}) + \mathrm{sin}^{2} \theta \cdot J_{0}(k r_{2}) \right)^{2} d \theta dx
        \ ,
        \label{e3_J0_W_int}
    \end{equation}

    \noindent
    where N is a normalization constant.


% ----------------------------------------
% Subsection 3.4: Expected results
% ----------------------------------------
\subsection{\label{sub:Expected rsults}Expected results}


    Let us reconfirm ``the electron phase'' in this research. As seen in subsection \ref{sub:Review_the_Previous_Study}, an electron has an internal structure. Each of the two spinors radiate its energy harmonically through time-dependent oscillation. Therefore, one electron obtain the phase of oscillation.

    Figure \ref{f_Biprism_2} illustrates how two bare electrons pass through the double slit. The summation of the amplitude of the two bare electrons is set to the value 100. Numbers in red letters indicate the distribution ration of the amplitude of both bare electrons. A red number of 100 means that all the energy possessed by the bare electron passes through either the left or right slit with 100 percent amplitude value. This means that an electron passes only through the slit on each side with whole energies. As noted later, no interference fringes are observed when passing through the slit at this phase.

        \begin{figure}[h!]
            \centering 
            \includegraphics[scale=0.4]{f_6_sum_edit.png}
                \caption{The summation of six sample patterns represented in Fig. \ref{f_6_edit}. The value of $P_{vis}$ is not equal to 1 as noted in Fig. \ref{f_Px=0}. (Yellow dashed circles)}
            \label{f_6_sum_edit}
        \end{figure}

    However, when bare electrons pass through the double slit at a ratio of 50 : 50, the situation is the same as a plane wave passed through. In other words, waves of the same amplitude traveled the double slit. Therefore, if the electrons are ejected from the electron gun in phase, a clear interference fringe with $P_{vis} = 0$ would be observed.
    
    In phase $0\pi$ (Fig. \ref{f_Biprism_2}), spinor 1 $(T_{\mathrm{e1}})$ occupies all the energy of the system in the electron. In phase $\theta_{\mathrm{phase \ at \ Biprism}} = 3/2\pi$, the ratio of amplitudes is the same as in phase $1/2\pi$. 

    Applying the electron model to the double slit experiment, interference fringes would not appear at a specific phase. On the specific phases, each of $A_{\mathrm{spinor1}}$ or $A_{\mathrm{spinor2}}$ equal to zero.


% ----------------------------------------
% Subsection 3.5: Simulation results
% ----------------------------------------
\subsection{\label{sub:Simulation results}Simulation results}


    In this subsection, we will demonstrate by simulation what the fringes look like for the 6 sampled phases at equally angled radians (see Table. \ref{table_1}). The simulation results are shown in Figs. \ref{f_6_edit} and \ref{f_6_sum_edit}.

    The difference of the interference fringes generated when two different amplitudes enter the double slit is clarified. Let us verify by computer simulation what kind of interference fringes appear depending on the phase of the electron. The program is described by Python and its source code is noted in the Appendix.

    Fig. \ref{f_6_edit} shows how the interference fringes change depending on the phase of the electrons. Fig. \ref{f_6_sum_edit} shows the total value of the sampled six patterns.
    
    The intensity was not zero value on the troughs of the fringes. This indicates that interference fringes with perfect contrast could not occur in electron experiments in principle.
    
    These results enforce the basis that the visibility in the two-slit experiment with spinor particles could not obtain value zero in theory.


    \begin{center}
        \begin{threeparttable}
            \caption{Six radians as samples.}
            \label{table_1}
            \tabcolsep = 6mm
            \renewcommand\arraystretch{1.2}
            \begin{tabular}{ccc} \hline
                radian & $A_{\mathrm{spinor1}}:\mathrm{cos^{2}} \theta$ & $A_{\mathrm{spinor2}}:\mathrm{sin^{2}} \theta$ \\ \hline
                0.785 & 0.500 & 0.500 \\
                0.628 & 0.655 & 0.345 \\
                0.471 & 0.794 & 0.206 \\
                0.314 & 0.905 & 0.095 \\
                0.157 & 0.976 & 0.024 \\
                0.000 & 1.000 & 0.000
                \\ \hline
            \end{tabular}
        \end{threeparttable}

    \end{center}

    \begin{center}
        \begin{threeparttable}
            \caption{Order of each particles. ($E \propto A^{2}$)}
            \label{table_2}
            \tabcolsep = 3mm
            \renewcommand\arraystretch{1.8}
            \begin{tabular}{cccc}
                \hline
                    & Spinor1 & Spinor2 & Vector \\
                \hline
                Energy & $\cos^{4} (\theta /2)$ & $\sin^{4} (\theta /2)$ & $\sin^{2} \theta$ \\
                Order & 4 & 4 & 2 \\
                \hline
                Amplitude & $\cos^{2} (\theta /2)$ & $\sin^{2} (\theta /2)$ & $\sin \theta$ \\
                Order & 2 & 2 & 1 \\
                \hline
            \end{tabular}
        \end{threeparttable}
    \end{center}


% = = * * = =  * * = = * * = = * * = = * * = = * * = = * * = = * * 
% ========================================
% Section 4: DISCUSSION
% ========================================
\section{\label{sec:DISCUSSION}DISCUSSION}


% ----------------------------------------
% Subsection 4.1: Amplitude order of spinor and vector particles
% ----------------------------------------
\subsection{\label{sub:Amplitude order of spinor and vector particles}Amplitude order of spinor and vector particles}


    Generally, the displacement of a simple harmonic oscillation takes the position of A at the upper limit and --A at the lower limit with reference to the origin. Trigonometric functions are useful for expressing the domain of displacement, for example:

    \begin{equation}
        x = A \sin (\omega_{\mathrm{0}} t + \delta), \ \ \ \ \  \omega_{\mathrm{0}} = \sqrt{\frac{k}{m}}
        \label{e_Simple_H_Osc}
    \end{equation}
    
    \noindent
    where $\omega_{\mathrm{0}}$ is angular frequency, $\delta$ is initial phase, $k$ is proportional constant, and $m$ is mass. The order of the trigonometric function in eqn \ref{e_Simple_H_Osc} is primary. Therefore, the domain, $x$, takes values from $A$ to $-A$, because $ -1 \leq \sin \theta \leq +1$.

    Referring to eqn \ref{e_E=E1+E2+gamma}, the first and the second term on the right side which are represented to the spinor particle are a forth order of trigonometric function, $\cos^{4} (\theta / 2)$ and $\sin^{4} (\theta / 2)$. On the other hand, the third term which are represented to the photon are second order, $\sin^{2} \theta$.

    As we have seen so far, bare electrons, which are Thermal Points, must have an energy fluctuation range of 0 to $A$. Because it does not become negative TPE value. Therefore, it may be reasonable to use a second order rather than a first order of trigonometric functions.
    
    Refer to eqn \ref{e_S1_S2_P} to reinforce this idea. Spinor particles taking an angle ($\theta$ / 2) are represented by the fourth order of trigonometric functions, $\sin^{4} (\theta / 2)$. On the other hand, a vector particle taking an angle $\theta$ is represented by the second order of trigonometric functions, $\sin^{2} \theta$.

    This definition yields the following hierarchy. That is, since the energy is directly proportional to the square of the amplitude of the oscillators, i.e. ($E \propto A^{2}$). The energy of spinor particles are expressed by the fourth order of trigonometric functions, and the amplitude of the spinor particles are expressed by the second order of trigonometric functions.
    
    The sum of the three terms represents the energy value, despite the fact that the trigonometric functions have different dimensions. The total value of both constitutes are the unit of energy, even though the trigonometric functions have different dimensions. Table. \ref{table_2} summarizes the relationship of these particles' orders. It is the relationship between energy and amplitude considered in this research. Validity of assign both $\cos^{2} \theta$ and $\sin^{2} \theta$ as weighting variables (eqn \ref{e3_Amplitude_AB} ) would need further study.


% ----------------------------------------
% Subsection 4.2: Prediction of the difference between experimental results for photons and electrons
% ----------------------------------------
\subsection{\label{sub:Prediction of the difference between experimental results for photons and electrons}Prediction of the difference between experimental results for photons and electrons}


    While the double-slit experiment of photons observes the interference fringe of $P_{vis} = 1$. How should we consider that the double-slit experiment of the electron, C$_{60}$ and helium atoms build-up the interference fringe of $P_{vis} \neq 1$ (see Fig. \ref{f_Px=0})?

    This subsection highlights the difference between the visibility of interference fringes occurs between photons and electrons, which depend on the structure of the electronic model from the previous study. This is because there are a pair of bare electrons in the electron, and these bare electrons constitute a model that vibrates by emitting and absorbing energy. This phenomenon occurs in spinor particles but not in vector particles. This is because there is no thermal radiation with TPE in the vector particle.

    Till date, the wave generated by the particle is regarded as a plane wave until the particles reached the double-slit. It is assumed that the wave takes equal amplitude values when passing through the double-slit. In this study, this phenomenon would be considered to be that two spinors, which are virtual spinners and bare electrons, split into two directions and pass through each of the double-slit.
    

% = = * * = =  * * = = * * = = * * = = * * = = * * = = * * = = * * 
% ========================================
% Section 5: CONCLUSION
% ========================================
\section{\label{sec:CONCLUSION}CONCLUSION}


    In this study, we can explain why the interference fringes produced by electrons passing through the double-slit do not have $(Visibility) = 1$. According to the previous study \cite{Hanamura}, we could get a clue which slit the electron passes through and the reason why no interference fringes appear when the electron is measured. That is, the experiment to determine which slit the electrons passed could be measured only when the phase of the electron is an integer multiple of $\pi$. In other words, in the phase $\theta = n \pi (n = 0,1,2,...)$, electrons pass through either the left or right slit. This is the case when $W_{1} = \mathrm{cos}^{2} \theta = 1$ or $W_{2} = \mathrm{sin}^{2} \theta = 1$.
    
    Furthermore, this study suggested that the interference fringes obtained when photons or electrons pass through the double-slit could be different. In the experiments of photons, very clear interference fringes appear. This may be considered these experiments obtain $(Visibility) = 0$. The difference between the two kinds of particles may be due to the fact that the electrons are fermions, on the other hand, the photons are bosons.
    
    This is because the interference fringes generated by photons would be equal intensity at the point of both slits, and it is reasonable that the conventional plane wave could be assumed.
    
    However, in the double-slit experiment using electrons, what this research would like to emphasize is that the two bare electrons' intensities differ depending on the electron's phase when they are passing through the double-slit.
    
    As proof that reinforces this issue, the C$_\mathrm{60}$ molecules double slit experiment gives us an example \cite{Arndt}. In the C$_\mathrm{60}$ experiment, striped patterns are observed on the left and right around the origin. And there is no clear dark line between stripes which indicated $P_{vis}$ is not zero value. In the double-slit experiment using Helium atomic beams \cite{Carnal}, interference fringes without clear dark lines observed as well. The discussion in this study should be further explored that differences in visibility observed in the double-slit experiments with photons and electrons.
    
    One way of verifying the validity of this electron model is to  control the electron phase in which the test electron emitted from the electron gun passes through the double-slit. If it becomes possible to control the phase of electrons on which passing through the biprism, it will be possible to change the visibility of the interference fringes intentionally. In the above case, even if one side of the double-slit is not closed, i.e., both slits are open, the double slit experiments could yield without interference fringes.
    
    If technology improves, it will be possible to verify the difference in the characteristics of the interference fringes build-up by electrons and photons. That is the technology which can adjust arbitrary phase of electrons passing through the double slit. Furthermore, it will be possible to perform double-slit experiments of electrons without interference or with complete contrast fringes.


% = = * * = =  * * = = * * = = * * = = * * = = * * = = * * = = * * 
% ========================================
% Bibliography
% ========================================
\begin{thebibliography}{99}
    \bibitem{Tonomura} A. Tonomura, J. Endo, T. Matsuda, T. Kawasaki and H. Ezawa. Demonstration of single-electron buildup of an interference pattern. $Amer. \ J. \ Phys. \ \bm{57}$ No.2 117--120 (1989). %

    \bibitem{Hanamura} S. Hanamura, A Model of an Electron Including Two Perfect Black Bodies, Zenodo (2018), \href{https://doi.org/10.5281/zenodo.16759284}{doi:10.5281/zenodo.16759284}. [Originally: \href{http://vixra.org/abs/1811.0312}{viXra:1811.0312}]

    \bibitem{Saiensu-sha} Hiroshi Ezawa, The Double-slit experiment (To the memory of Akira Tonomura), Mathematical Sciences; Saiensu-sha Co., Ltd.: Tokyo, No.635, (May 2016) pp. 9--14 (written in Japanese).

    \bibitem{Omar} Omar S. Mag\~{n}a-Loaiza, $et \ al.$, Exotic Looped Trajectories of Photons in Three-Slit Interference. $nature \  communications$. DOI: 10.1038 (2016).

    \bibitem{Arndt} Arndt, Markus, $et \ al.$, Wave-particle duality of C$_{60}$ molecules. $Nature$. 401. 680-682. (1999).
    
    \bibitem{Carnal} Carnal, O \& Mlynek, Jo. Young’s double-slit experiment with atoms: A simple atom interferometer. $Physical \ review \ letters. \ \bm{66}$. 2689-2692. (1991). 


    \
    
        \begin{flushright}
            (This paper has been unedited.)
        \end{flushright} 
    

\end{thebibliography}

\clearpage

\onecolumngrid


% = = * * = =  * * = = * * = = * * = = * * = = * * = = * * = = * * 
% ========================================
% Section 6: Appendix
% ========================================
\section{\label{sec:Appendix}Appendix}


    \begin{figure*}[h!]
        \centering
        \includegraphics[scale=0.6]{I_listing1.png}
        \caption{The source program of Fig. \ref{f_6_edit}. The interference fringes were simulated in lines from 16 to 22.}
    \label{I_listing1}
    \end{figure*}

    \begin{figure*}[h!]
        \centering
        \includegraphics[scale=0.6]{I_listing2.png}
        \caption{The source program of Fig. \ref{f_6_sum_edit}. The value of the vertical axis (Intensity) was omitted near the origin. This is because the simulation values interfere near the origin. This situation was described in the sixth line, where the plot range of the vertical axis was from 0.00001 to 0.00016.}
        \label{I_listing2}
    \end{figure*}

    \begin{figure*}[h!]
        \centering
        \includegraphics[scale=0.6]{I_f_6_edit_2.png}
        \caption{The result of the simulation which plotted the horizontal axis from the origin. Since the peak value of the intensity was recorded strongly around the origin, it was difficult to find interference fringes. This was due to the characteristics of the zero-order Bessel function. Therefore, in this research, intensity was collected at a place away from the origin.}
        \label{I_f_6_edit_2}
    \end{figure*}

    \begin{figure*}[h!]
        \centering
        \includegraphics[scale=0.6]{I_listing3.png}
        \caption{The source program of Fig. \ref{I_f_6_edit_2}. This source was changed lines 16 and 17 of Fig. \ref{f_6_edit}. In this research, the value within the blue circle was set to 4000 instead of 0. This change made it easy to obtain interference fringes because the intensity was clearly observed away from the origin.}
        \label{I_listing3}
    \end{figure*}
    

% ========================================
% End of Document
% ========================================
\end{document}
